% !TEX root = ../main.tex

\chapter{App Development and Testing}

This chapter reports what the MNEME front end and back end actually implement and what the retained test evidence shows. Implementation claims are bounded by two inspected repository revisions: the integrated repository state \texttt{b945d3c}, which contains the backend service, the AI service layer, the behavior baseline, the citation-graph platform, the integrated Android client, and their test suites; and the AI evaluation branch tip \texttt{e11e385}, which adds the graded evaluation harness, the pinned fixture corpus, and the live evaluation command. A capability that appears in neither revision is a design statement and belongs to Chapter~3, not here.

Throughout the chapter, software tests are interpreted narrowly: a passing test establishes the specified behavior it checks, in the environment in which it ran, and nothing more. Measured latency distributions are reported with their sample sizes and collection conditions, estimated costs are labeled as estimates, and quality properties that would require a live graded run against a deployed model provider are stated as limitations rather than anticipated results. This discipline connects the chapter to the technical challenges of Chapter~1: the mechanisms below are the implemented responses to representation continuity, evidence continuity, interest continuity, and interaction continuity, and the testing sections report which of those responses currently carry executable or measured evidence.

\section{Front-end Development}

The Android client is a Kotlin application built on Jetpack Compose with Material~3 components. Its role in the system is deliberately narrow: it renders server-produced research artifacts --- briefings, paper records, summaries, grounded answers, and a bounded citation graph --- while keeping the user's reading state recoverable when the network, the backend, or the device lifecycle interrupts the session. That narrowness is the client-side response to the interaction-continuity challenge: continuity is achieved not by making the client smart, but by making every visible state explicit about where its content came from and what the user can do next.

\subsection{Design Pattern}

The client follows the Model--View--ViewModel pattern with a repository layer between ViewModels and data sources. Compose screens are pure functions of observable UI state; ViewModels own that state and reduce it from repository results; repositories mediate between Retrofit/OkHttp network services, Room persistence, and WorkManager background execution. The pattern was chosen for two reasons that matter to this thesis rather than for framework fashion. First, it makes state transitions unit-testable without an emulator: the controlled state measurements reported under Testing Results drive ViewModel reducers directly with deterministic fixture inputs. Second, it gives the offline and degraded paths a single home --- the repository decides whether a screen receives live content, cached content, or an error, and the decision is visible in the state object rather than buried in view code.

A distinctive element of the state model is its disclosure vocabulary. Every content-bearing state distinguishes loading, live, cached, controlled fixture, empty, fallback, and retryable error paths where applicable, and the distinction is rendered to the user rather than kept internal. A briefing restored from the local cache is labeled as cached with its freshness; a citation graph degraded to its bounded local fallback says so; controlled fixture content used in development builds is marked as such. This design choice follows directly from the evidence discipline of the thesis: if the manuscript must never promote cached or fixture content into live evidence, the running application should not do so either.

The repository layer also carries the client's consistency obligations. Briefings and paper records are retained locally with retention metadata; behavioral events are queued durably with client-generated UUID identity and retry bookkeeping, so that a process death between capture and upload loses nothing and a repeated upload can be recognized by the server as a duplicate rather than a new signal. These mechanisms are exercised by the event-synchronization results reported later in this chapter.

\subsection{UI Controls}

The home surface is the briefing list, rendered as Compose cards that pair each recommended paper with the match information and human-readable reason produced by the backend scorer, and one dominant open action. The paper detail screen presents the structured summary with its source-linked elements, and the question-answering surface accepts a free-form question and renders the grounded answer together with its citation markers, so that a reader can move from a claim to the evidence chunk that supports it. Verification labels sit adjacent to the claims they qualify rather than in a separate legend; this placement was a direct consequence of the formative usability observations reported in Chapter~3.

The citation graph is the one control that leaves the Compose world. It is rendered by a local D3 bundle inside a WebView, with no network fetch of the visualization library, while node selection is mirrored back into native state through a JavaScript bridge. The native side therefore remains the single source of truth for which paper is selected, which allows the persistent ``open paper'' action, the unambiguous selected-node style, and the selection restoration behavior verified under UI Testing. Graph size is bounded by the client request contract at fifty nodes, which keeps both the rendering cost and the user's orientation manageable on a phone screen.

Interest management uses standard Material list and edit controls, reachable from the briefing surface, through which explicit topics can be added, edited, and removed. Status and error surfaces reuse a small set of shared components so that the disclosure vocabulary described above looks identical wherever it appears.

\subsection{Navigation Controllers}

Navigation uses Compose Navigation with type-safe routes. The primary flow follows the research loop the system is designed around: briefing to paper detail, paper detail to question answering, paper detail or search to the citation graph, and briefing to interest and preference settings. Each route carries typed arguments such as paper identity and briefing identity rather than parsed strings, so a navigation regression is a compile-time failure rather than a runtime surprise.

Two navigation behaviors carry evidence in this chapter. First, returning from a paper opened out of the graph restores the previous graph selection rather than resetting the view; this is part of the implemented graph and app-flow test suite. Second, the back stack is arranged so that background refresh and notification entry points land on the briefing surface with state restored from the local cache when the process was killed, which is what makes the cold-start and resume measurements reported under Performance Testing meaningful user-visible quantities rather than synthetic numbers.

\subsection{Other Aspects}

Persistence is provided by Room with explicit migrations through schema version~3; briefing and paper caches carry offline metadata and retention policies. Background work is scheduled through WorkManager, which owns periodic refresh and recovers interrupted work after process death; the durable event queue described above is drained by this machinery, and a reserved event returns to the pending state when its worker is interrupted. Local notification primitives for new-briefing delivery are implemented, and first-run onboarding derives an initial arXiv category from a seed paper and prepares an initial set of papers through the live backend pipeline before the first briefing is shown.

The client's scope boundary is stated here once so that the testing sections do not have to hedge repeatedly. Cache-backed content states, event queueing and recovery, graph interaction, and the live client-to-backend event loop are present in code at the stated revisions and carry the evidence reported in this chapter. The complete notification permission and delivery flow across Android versions, and broader device coverage beyond the measured emulator configuration, remain outside the retained evidence and are treated as limitations rather than claims.

\section{Back-end Development}

The backend is a Python/FastAPI service with PostgreSQL for durable storage, \texttt{pgvector} for vector search, Redis for caching, budget accounting, and job queues, and ARQ workers for background processing. Its data model separates a paper's logical identity from its observed arXiv revisions: metadata ingestion registers papers, while long-running preparation --- parsing, section-aware chunking, summarization, and embedding --- runs in workers whose job state survives process restart. Every derived artifact binds to a specific paper revision, which is the structural foundation for the representation-continuity mechanisms described next.

Representation continuity is enforced by determinism and identity rather than by convention. Chunking is deterministic: the same parsed sections always produce the same chunk contents and content hashes, which makes the embedding stage idempotent and re-chunking comparable across runs. The frozen schema stores each embedding in a 1536-dimension vector column paired with the identity of the model that produced it, under a check constraint that either both or neither are set, so a model swap cannot leave orphan vectors that silently poison similarity search. Retrieval is revision-scoped by construction, and an API test asserts that question answering over one revision cannot return chunks of another; a new arXiv revision therefore produces new artifacts rather than contaminating old ones.

Evidence continuity is implemented as a pipeline that is honest about each of its stages. Retrieval embeds the question and performs dense approximate-nearest-neighbor search over one revision's chunks, returning the top eight by cosine similarity plus up to two early-document context anchors that preserve overview context; dense membership and anchor membership are tracked independently, because an early section can legitimately be both. Reranking is a transparent lexical blend rather than a learned cross-encoder: the final score weights vector similarity at 0.7 and the fraction of the question's content words present in the chunk at 0.3, ties break on original retrieval order, and the top four chunks become evidence, with anchors outside that set appended rather than displacing question-specific evidence. Generation is refused before any model call when no evidence reaches a minimum score of 0.25, the model may itself decline through a designated insufficient-evidence marker, and both refusal paths return the same fixed refusal answer, so the client never has to distinguish them. After generation, citation verification checks every bracketed marker against the evidence list --- markers outside the list are dropped rather than repaired --- and compares each cited chunk with the local claim sentences that carry its marker; a citation is source-matched when at least 30\% of the claim's content words appear in the cited chunk, and answers are labeled matched, partially matched, or insufficient. This is retrieved-source matching, a lexical construct, and the thesis does not present it as semantic entailment or scientific truth.

Interest continuity is served by a recommendation scorer that composes three signals: explicit topic match with weight 0.45, cosine similarity between the user's learned behavior embedding and a paper's mean chunk embedding with weight 0.35, and publication recency with a one-week half-life at weight 0.2. Missing signals redistribute their weight over the available ones, so a cold-start user still receives a meaningful ranking, and the ranking is deterministic with a stable tie-break. Each recommendation carries human-readable reasons derived from its dominant signals rather than free-form generated text, which is what the client renders beside each briefing entry. The behavior embedding itself is produced by the behavior model owned by the data platform; this chapter's claims concern the composition, not the internals of that model.

Cross-cutting the pipeline are provider routing, caching, and cost control. Each AI task routes to a configured provider; Anthropic, OpenAI, and DeepSeek adapters are implemented, providers are inferred from model identifiers, and a misconfigured route fails at service startup rather than on the first user request. Completions are cached in Redis with a 24-hour time-to-live for question answering and seven days for summaries, with the prompt version participating in every cache key so that a template change invalidates cleanly. A shared daily budget guard refuses provider calls once the configured daily cap (five US dollars by default) is spent, and exhaustion surfaces to the client as an explicit rate-limit response rather than a silent failure. Every completion persists provider, model snapshot, prompt version, input hash, token counts, estimated cost, and latency, so any live run is auditable from the database without additional instrumentation.

Because evidence continuity is a headline challenge, the backend also implements the instrumentation needed to measure it, not only to provide it. A graded evaluation harness drives the full retrieval and grounded-answer pipeline over a frozen fixture corpus: fifteen manually curated cases across seven arXiv papers, twelve answerable and three whose correct behavior is refusal. Each fixture pins the exact arXiv revision it was authored against, and the live evaluation command resolves that exact revision --- never ``latest'' --- and requires every chunk of it to carry an embedding produced by the configured embedding model before a case may run; anything else is reported as a skip with explicit counts. Retrieval quality is graded on dense results only, so an appended context anchor cannot convert a retrieval miss into a hit. Every report embeds a run-configuration block recording the provider, model, prompt version, embedding backend and model, retrieval and reranking settings, evidence threshold, and cache posture, and separates the served completion's generation metadata from the cost and latency the run actually incurred, so a cached run cannot masquerade as live spend. A run that evaluates zero fixtures exits with a nonzero status so that automation cannot mistake an empty evaluation for a healthy one. This infrastructure exists at revision \texttt{e11e385}; the graded quality numbers it will produce require a live run against ingested fixture papers and are therefore not reported in this thesis.

Taken together, these mechanisms are the implemented answers to the Chapter~1 challenges: revision-bound artifacts and idempotent preparation for representation continuity; bounded retrieval, refusal, and citation verification for evidence continuity; the composed, explainable scorer for interest continuity; and the job, event, and briefing boundaries that the mobile client's recoverable loop depends on for interaction continuity. The remainder of the chapter reports the evidence that currently backs them.

\section{Testing Results}

The project distinguishes four evidence layers, and the distinction governs what each result below is allowed to mean. Unit tests check deterministic transformations and boundary validation; repository and API tests check transactions, identity, idempotency, and controlled responses; integration and device tests check that client and server agree and that visible state follows backend outcomes; and research evaluations measure quality or usefulness on frozen data or with participants. Passing a lower layer is necessary but never substituted for a higher one. The results reported here occupy the first three layers, together with bounded device measurements; the graded answer-quality evaluation belongs to the fourth layer and, as stated above, its execution remains outstanding.

\subsection{Testing Tool}

The backend test suite runs under \texttt{pytest} in a \texttt{uv}-managed Python environment, with \texttt{ruff} for lint and formatting and \texttt{pyrefly} for type checking as static gates. Provider-dependent code is tested against deterministic fake LLM and embedding providers and an in-memory Redis stand-in, so no software test in this thesis depends on, or is presented as evidence about, a live model. At the evaluation branch tip \texttt{e11e385} the suite reports 463 passed and 7 skipped, the skips being PostgreSQL-only integration tests that require a configured database; the integrated revision \texttt{b945d3c} carries its own passing suite under the same tooling in continuous integration.

The Android suite uses JUnit-based unit tests for ViewModel and repository logic, Compose UI tests for screen behavior, and Gradle-orchestrated static analysis with ktlint, detekt, and Android lint. Device measurements were collected on one Android~14 (API~34) Google emulator with four virtual processors and a 192~MiB application heap. For every measured group, warm-up iterations were run and discarded before retained samples were collected; reported medians are sample medians and 95th percentiles use the nearest-rank definition. This single-configuration environment is a stated boundary of every device number below: it supports controlled comparison across conditions, not an estimate of physical-device variability.

\subsection{UI Testing}

The client's visible-state machinery was exercised in controlled trials that deliver deterministic fixture results to the ViewModel layer and verify the resulting state. Five content-state conditions --- a live seed transition, a cached warm start, a cached offline disclosure, a cached server-failure disclosure, and an invalid-cache error --- each completed 30 of 30 trials with the expected state and disclosure label. The three citation-graph disclosure states (ready, fallback, and empty) likewise each completed 30 of 30 trials. These results establish that the disclosure vocabulary described under Design Pattern is actually reachable and correctly labeled under its specified conditions; because the transitions are driven by fixtures, they say nothing about network behavior or rendering performance, which are treated separately below.

Interactive graph behavior was verified by the graph and app-flow suite, which passed 11 of 11 tests, including node selection through the WebView bridge, navigation to a selected paper through the persistent open action, and restoration of the previous selection after returning from paper detail. Together with the state trials, this establishes the specified presentation and navigation behavior of the graph surface. It does not measure graph usefulness, discoverability, or layout quality; the formative five-participant observations that motivated the graph's redesign, and the redesign itself, are reported in Chapter~3 as design evidence rather than repeated here.

\subsection{Acceptance Testing for Features}

The feature-level acceptance evidence maps the acceptance criteria of Chapter~3 to executable checks at the stated revisions. For the literature-representation features, the decisive check is revision isolation: an API test asserts that question answering scoped to one paper revision cannot retrieve chunks belonging to another, which is the executable form of the promise that a revision update cannot silently change what the user is reading. Idempotent re-preparation is covered by the deterministic chunking and content-hash tests in the backend suite.

For the assistance features, acceptance rests on the refusal and verification behavior. Tests cover refusal both below the evidence threshold and through the model-side insufficient-evidence marker, citation-marker validation against the evidence list, and the source-match labeling of answers. The budget boundary is covered on the public question route: integration tests assert that exhaustion of the daily budget surfaces as the specified rate-limit response for both the completion guard and the embedding guard, rather than as a generic failure. The evaluation command itself carries contract tests for its fixture pinning, its readiness checks with explicit skip reasons, its provenance snapshot, and its nonzero exit on an empty run.

For the interaction loop, the acceptance evidence is the event-synchronization behavior measured end to end. In controlled trials isolating the client queue from network variability, 30 of 30 trials each succeeded for the durable queue write, the injected retryable failure, the subsequent retry success, and the duplicate replay. A live trace then exercised the same client coordinator against a running backend in five repetitions: each repetition queued three new events, recovered from an injected retryable failure, uploaded the batch, replayed the same identifiers, and read a briefing back through the standard client data path. Every live upload accepted three new identifiers, every replay was acknowledged as duplicate rather than re-ingested, and every readback carried a live-backend origin. The returned digest contained zero papers in all five repetitions because the evaluation database held no paper inside the active fourteen-day briefing window, and the preference-model version remained~1 throughout; the live loop therefore demonstrates reachability and duplicate-aware ingestion of the cross-layer path, and deliberately does not demonstrate recommendation content or adaptation. Asynchronous job visibility was verified by the three-stage paper-preparation polling flow, which completed in 10 of 10 trials.

\subsection{Performance Testing}

Client-side performance was measured on the emulator configuration described under Testing Tool. Figure~\ref{fig:mobile-reliability} and Table~\ref{tab:mobile-reliability} report the retained distributions. All 20 cold process starts and all 20 task foreground/resume operations returned a wait-time record: cold start had a median of 2{,}773~ms with a 95th percentile of 4{,}444~ms, and foreground/resume had a median of 133.5~ms with a 95th percentile of 445~ms. The startup runner records activity wait time and does not inspect rendered content; Android classified nine resume operations as hot starts and returned an unknown classification for eleven, so the resume row is reported as one distribution rather than as twenty confirmed hot starts. The five controlled content-state conditions completed with sub-millisecond medians, which is expected --- they time the ViewModel transition after a deterministic fixture result is delivered, and must not be read as network or display latency. The three-stage job polling flow's median of 3{,}008~ms is consistent with its one-second client polling interval across three processing stages.

\begin{figure}[htbp]
\centering
\includegraphics[width=0.92\textwidth]{figures/mneme_mobile_reliability.pdf}
\caption{Android client latency and reliability measurements on the single emulator configuration. Separate panels keep process/task timing, controlled content-state transitions, and event operations distinct, because the three groups measure different quantities and cannot be pooled into one timing family.}
\label{fig:mobile-reliability}
\end{figure}

\begin{table}[htbp]
\centering
\caption{Android startup, controlled state, and polling results. Latencies are in milliseconds; p95 uses the nearest-rank definition.}
\label{tab:mobile-reliability}
\begin{tabular}{p{5.1cm}rrrr}
\toprule
Condition & $n$ & Successful & Median & p95 \\
\midrule
Cold process start & 20 & 20 & 2{,}773 & 4{,}444 \\
Task foreground/resume & 20 & 20 & 133.5 & 445 \\
\addlinespace
Live seed transition (controlled) & 30 & 30 & 0.124 & 0.293 \\
Cached warm start (controlled) & 30 & 30 & 0.227 & 0.376 \\
Cached offline disclosure (controlled) & 30 & 30 & 0.227 & 0.439 \\
Cached server-failure disclosure (controlled) & 30 & 30 & 0.231 & 0.362 \\
Invalid-cache error (controlled) & 30 & 30 & 0.203 & 0.635 \\
\addlinespace
Three-stage job polling (controlled) & 10 & 10 & 3{,}008.418 & 3{,}064.345 \\
\bottomrule
\end{tabular}
\end{table}

Graph presentation cost was measured at 1, 12, 25, and 50 nodes --- the upper bound matching the client request contract --- in both a newly created WebView and an update to an existing WebView. Figure~\ref{fig:graph-latency} shows the distributions and Table~\ref{tab:graph-latency} their summary statistics. Every ready and selection sample succeeded at every size. Median DOM-ready latency ranged from 139.749 to 454.149~ms in a new WebView and from 83.001 to 111.814~ms in an existing one, and median selection-callback latency stayed between 32.279 and 49.098~ms across all sizes, which supports the design decision to keep one WebView alive and update it rather than recreate it per navigation. The medians are not monotone in node count and several 95th percentiles are wide; this reflects emulator scheduling and WebView initialization variability, and the experiment accordingly supports no linear scaling model. Render timing ends at the local graph page's DOM-readiness report, while the force-directed layout may still be converging, so these numbers bound presentation readiness, not visual settlement.

\begin{figure}[htbp]
\centering
\includegraphics[width=0.92\textwidth]{figures/mneme_graph_latency.pdf}
\caption{Citation-graph presentation and selection latency by node count. Render timing ends when the local graph page reports DOM readiness; it does not measure convergence of the force-directed layout.}
\label{fig:graph-latency}
\end{figure}

\begin{table}[htbp]
\centering
\caption{Graph DOM-ready and native selection-callback latency. Each cell is median~/~p95 in milliseconds.}
\label{tab:graph-latency}
\begin{tabular}{p{1.1cm}p{3.0cm}p{3.0cm}p{3.0cm}p{3.0cm}}
\toprule
Nodes & New WebView ready & Existing WebView ready & New WebView selection & Existing WebView selection \\
\midrule
1  & 139.749 / 298.019 & 83.001 / 150.840 & 32.764 / 46.022  & 32.279 / 103.680 \\
12 & 454.149 / 913.746 & 95.590 / 331.895 & 35.971 / 392.792 & 32.968 / 208.816 \\
25 & 315.039 / 734.190 & 97.767 / 361.523 & 48.675 / 568.963 & 32.759 / 128.642 \\
50 & 306.793 / 928.176 & 111.814 / 513.785 & 49.098 / 353.710 & 37.860 / 163.644 \\
\bottomrule
\end{tabular}
\end{table}

Event-loop performance completes the client-side picture. Table~\ref{tab:event-sync} reports the controlled stages over 30 trials each and the live-backend stages over five repetitions; the two groups are kept separate because only the latter includes network and server time. Controlled queue writes carried a median cost of about 10~ms, and the live upload-after-retry stage a median of 33~ms with a 95th percentile of 188~ms, so the durability and duplicate-safety mechanisms described earlier operate at costs well below user-perceptible thresholds in the measured environment.

\begin{table}[htbp]
\centering
\caption{Client event-synchronization latency, controlled stages ($n=30$) and live-backend stages ($n=5$). Latencies are in milliseconds; all trials in every stage succeeded.}
\label{tab:event-sync}
\begin{tabular}{p{6.2cm}rrr}
\toprule
Stage & $n$ & Median & p95 \\
\midrule
Durable queue write (controlled) & 30 & 10.287 & 32.261 \\
Retryable HTTP failure (controlled) & 30 & 15.613 & 80.332 \\
Retry success (controlled) & 30 & 54.999 & 91.794 \\
Duplicate replay (controlled) & 30 & 18.950 & 50.222 \\
\addlinespace
Queue three client events (live) & 5 & 13.850 & 19.061 \\
Injected failure before retry (live) & 5 & 8.028 & 8.616 \\
Upload after retry (live) & 5 & 33.049 & 188.478 \\
Replay the same identifiers (live) & 5 & 24.210 & 125.869 \\
Read back a live briefing (live) & 5 & 26.776 & 135.906 \\
\bottomrule
\end{tabular}
\end{table}

On the backend side, model cost is treated as a first-class performance dimension because the assistant runs continuously against paid providers. A static audit inventoried every provider call site --- the summarization worker, grounded question answering, query embedding, and batched chunk embedding; recommendation scoring makes no provider call because it reuses stored vectors --- and estimated unit economics from provider list prices and the deployed input and output caps. Under those caps, one paper summary on the default flagship route costs an estimated \$0.10, one grounded answer roughly \$0.028, and embedding a full paper roughly \$0.0003; against the default \$5 daily budget, a single bulk-ingestion day of about fifty summaries can exhaust the cap on the flagship tier, while a lighter tier fits roughly 250 summaries per day, and a full fifteen-case evaluation run is estimated at \$0.34 versus \$0.07 across the two tiers. These figures are estimates, not billed spend: the underlying price rows were re-verified against provider pricing pages and retained, with the cap assumptions, in the repository artifact \texttt{docs/evaluation/ai-cost-price-snapshot-2026-07-26.md}, whose companion script recomputes every figure and fails if prose and calculation diverge. The audit's actionable finding --- that default routing sends bulk summarization, the dominant and least quality-sensitive spend, to the most expensive model --- was recorded as a recommendation and deliberately not applied, because the rerouting decision is gated on a summary-quality comparison that has not yet been run.

The boundaries of this section's evidence are stated plainly. All device numbers come from one emulator configuration and do not estimate physical-device variability. Controlled state and event timings measure client-internal transitions, not network or frame latency, and graph timing ends at DOM readiness. Cost figures are estimates derived from a retained price snapshot. Finally, the graded answer-quality run --- retrieval hit rate, source-match rate, refusal accuracy, and measured per-answer latency and cost on the pinned fixture corpus --- has not been executed against a live provider, so this thesis reports the evaluation instrumentation as implemented and verified, and reports no answer-quality numbers; completing that run is future work discussed in Chapter~5.
